\chapter{The Polyakov Path Integral, Ghosts and the Superstring in Brief}\label{ch:ghosts}

\Cref{ch:quantum} found the critical dimension $D=26$ by demanding Lorentz invariance in
lightcone gauge, a calculation in which the number $26$ appears at the end of a long
commutator. This chapter answers a sharper question: \emph{where does $26$ come from, and is it
really a number of dimensions?} The answer uses the conformal field theory of \cref{ch:cft}.
The worldsheet metric is a gauge field; removing it from the path integral costs a Jacobian;
the Jacobian can be written as a new conformal field theory of anticommuting ``ghost'' fields;
and that theory has central charge $c=-26$. Consistency of the gauge symmetry then demands
that everything else on the worldsheet has $c=+26$. A student should care for two reasons.
First, the result is a \emph{budget}: any conformal field theory with the right central charge
is a legitimate string background, which is what makes compactification on circles, tori and
K3 surfaces (the subject of most of this book) possible at all. Second, the same argument,
repeated for the superstring, gives $D=10$ and the budget $c=6$ for a K3 surface that returns
in \cref{ch:stringsk3,ch:ellgenus}. We end with a short, descriptive account of the
superstring theories, at the level needed later, and with an exact statement of which parts
of the chapter's arithmetic have been checked in Lean.

\section{The problem: a gauge symmetry with an anomaly}\label{sec:ghosts:problem}

Recall from \cref{ch:classical} that the Polyakov action\index{Polyakov action} for $D$
scalar fields $X^\mu(\sigma)$ on a worldsheet with metric $g_{\alpha\beta}$ is, in Euclidean
signature,
\begin{equation}\label{eq:ghosts:poly}
  S_{\mathrm{Poly}}[X,g]=\frac{1}{4\pi\ap}\int \dd^2\sigma\,\sqrt{g}\;
  g^{\alpha\beta}\,\partial_\alpha X^\mu\,\partial_\beta X^\nu\,\delta_{\mu\nu},
\end{equation}
\source{Tong §5.1, ll.~6300--6308}. It has two local symmetries. \emph{Diffeomorphisms}
$\sigma^\alpha\to\sigma'^\alpha(\sigma)$ relabel the points of the worldsheet. \emph{Weyl
transformations}\index{Weyl transformation} $g_{\alpha\beta}\to\ee^{2\omega(\sigma)}g_{\alpha\beta}$
rescale the metric point by point and leave $X^\mu$ alone; they are a symmetry because
$\sqrt{g}\,g^{\alpha\beta}$ is unchanged in two dimensions
($\sqrt{g}\to\ee^{2\omega}\sqrt g$ and $g^{\alpha\beta}\to\ee^{-2\omega}g^{\alpha\beta}$).

Both are \emph{gauge} symmetries. A gauge symmetry is not a symmetry of nature but a
redundancy of the description: two configurations related by it are the same physical state.
A global symmetry may be broken by quantum effects and the theory survives; if a gauge
symmetry is broken, the redundancy no longer cancels and the unphysical degrees of freedom
(here, the timelike oscillators with negative norm met in \cref{ch:quantum}) no longer
decouple. A gauge anomaly is therefore fatal \source{Tong §5, ll.~6274--6290}.

Now recall two facts from \cref{ch:cft}. First, on a curved worldsheet with Ricci scalar $R$,
the trace of the stress tensor of a conformal field theory with central charge $c$ does not
vanish:
\begin{equation}\label{eq:ghosts:weylanomaly}
  \vev{T^\alpha{}_\alpha}=-\frac{c}{12}\,R
\end{equation}
\source{Tong §4.4.2, ll.~5094--5112}. Since $T^\alpha{}_\alpha$ is the response of the theory
to a Weyl rescaling (it is $g^{\alpha\beta}\,\delta S/\delta g^{\alpha\beta}$ up to
normalization), \cref{eq:ghosts:weylanomaly} says that the quantum theory is \emph{not} Weyl
invariant unless $c=0$. This is the Weyl anomaly\index{Weyl anomaly}. Second, every
non-trivial unitary conformal field theory has $c>0$; for instance $D$ free scalars have
$c=D$. On a fixed background this is harmless. In string theory, where we integrate over the
metric, it looks like a contradiction: we need $c=0$ and cannot have it
\source{Tong §5, ll.~6290--6299}.

The loophole is that we have not yet computed the path integral correctly. Removing the gauge
redundancy produces additional fields, and these fields are not unitary: they carry
\emph{negative} central charge. The condition $c=0$ applies to the total.

\section{The Faddeev--Popov method}\label{sec:ghosts:fp}

\subsection{What ``dividing by the volume'' means}

The quantum string is defined by the path integral\index{path integral!Polyakov}
\begin{equation}\label{eq:ghosts:Z}
  Z=\frac{1}{\mathrm{Vol}}\int \mathcal{D}g\,\mathcal{D}X\;\ee^{-S_{\mathrm{Poly}}[X,g]},
\end{equation}
where ``Vol'' is the volume of the group of diffeomorphisms and Weyl transformations
\source{Tong §5.1, ll.~6316--6340}. The integrand is constant along each \emph{gauge
orbit}\index{gauge orbit} (the set of all configurations obtained from one by gauge
transformations), so the integral over all metrics overcounts each physical configuration by
the volume of its orbit. We want to integrate once over each orbit. To do so we choose a
\emph{gauge slice}: a surface in field space that meets every orbit exactly once
(\cref{fig:ghosts:orbits}). Changing variables from ``all fields'' to ``position on the
slice'' and ``position along the orbit'' costs a Jacobian. The method of Faddeev and Popov
\index{Faddeev--Popov method} computes this Jacobian for any gauge symmetry
\source{Tong §5.1, ll.~6341--6353}.

\begin{figure}[ht]
\centering
\begin{tikzpicture}[scale=0.95]
  \foreach \r in {0.7,1.3,1.9,2.5} {\draw[thick] (0,0) circle (\r);}
  \draw[dashed,very thick,blue!60!black] (0,0) -- (3.1,0) node[right] {slice $y=0,\ x>0$};
  \foreach \r in {0.7,1.3,1.9,2.5} {\fill[blue!60!black] (\r,0) circle (1.6pt);}
  \draw[-{Stealth}] (1.9,0) ++(40:1.9) arc (40:70:1.9);
  \node at (2.0,2.05) {\small orbit};
  \fill (0,0) circle (1.2pt);
  \draw[dotted,thick] (-3.0,0) -- (0,0);
  \foreach \r in {0.7,1.3,1.9,2.5} {\draw[thick,fill=white] (-\r,0) circle (2pt);}
  \node[align=right,left] at (-3.05,0) {\small second crossings of $y=0$\\\small (Gribov copies)};
\end{tikzpicture}
\caption{Gauge orbits (solid circles) of the rotation group acting on the plane, and a gauge
slice (dashed half-line) that meets each orbit once. The full line $y=0$ would meet each orbit
twice. In string theory the plane is replaced by the space of worldsheet metrics, the rotations
by diffeomorphisms and Weyl transformations, and the slice by a choice of fiducial metric
$\hat g$.}
\label{fig:ghosts:orbits}
\end{figure}

\subsection{A warm-up with numbers}

\begin{example}[Faddeev--Popov in the plane]\label{ex:ghosts:plane}
Let the ``fields'' be a point $\vec x=(x,y)\in\R^2$, the ``action'' $S=x^2+y^2=r^2$, and the
gauge group the rotations $\vec x\to\vec x^{\,\theta}$ by an angle $\theta\in[0,2\pi)$, with
$\mathrm{Vol}=2\pi$. The quantity to compute is
\[
  Z=\frac{1}{2\pi}\int_{\R^2}\dd x\,\dd y\;\ee^{-(x^2+y^2)}=\frac{1}{2\pi}\cdot\pi=\frac12 .
\]
Now repeat the computation the Faddeev--Popov way. Choose the slice $y=0$, $x>0$. Define
$\Delta(\vec x)$ by
\[
  \Delta(\vec x)^{-1}=\int_0^{2\pi}\dd\theta\;\delta\bigl(y^\theta\bigr)\Big|_{x^\theta>0},
  \qquad y^\theta=r\sin(\theta+\theta_0),
\]
where $\theta_0$ is the polar angle of $\vec x$. Using
$\int\dd\theta\,\delta(f(\theta))=1/|f'(\theta_*)|$ at the zero $\theta_*=-\theta_0$ gives
$\Delta^{-1}=1/(r\cos 0)=1/r$, so $\Delta(\vec x)=r$; it depends only on the orbit, not on
the position along it. Insert
$1=\Delta(\vec x)\int\dd\theta\,\delta(y^\theta)$ into $Z$, change variables
$\vec x\to\vec x^{\,-\theta}$ (the measure and the action are rotation invariant), and the
$\theta$ integral factors out as $2\pi=\mathrm{Vol}$:
\[
  Z=\int\dd x\,\dd y\;\Delta(\vec x)\,\delta(y)\Big|_{x>0}\,\ee^{-(x^2+y^2)}
   =\int_0^\infty\dd x\;x\,\ee^{-x^2}=\frac12 .
\]
The two answers agree. The Jacobian $\Delta=r$ is nothing but the $r$ in
$\dd x\,\dd y=r\,\dd r\,\dd\theta$. Had we forgotten it we would have obtained
$\int_0^\infty\ee^{-x^2}\dd x=\sqrt\pi/2\approx0.886$, which is wrong. Had we used the whole
line $y=0$ as the slice we would have counted each orbit twice; such repeated intersections
are called \emph{Gribov copies}\index{Gribov copies} \source{Tong §5.1.1, ll.~6396--6401}.
\end{example}

\subsection{The same steps for the string}

Write $\zeta$ for a combined diffeomorphism and Weyl transformation, and $g^\zeta$ for the
transformed metric,
$g^\zeta_{\alpha\beta}(\sigma')=\ee^{2\omega(\sigma)}\,
\frac{\partial\sigma^\gamma}{\partial\sigma'^\alpha}\frac{\partial\sigma^\delta}{\partial\sigma'^\beta}\,
g_{\gamma\delta}(\sigma)$. A metric in two dimensions has three independent components, and
$\zeta$ contains three arbitrary functions ($\omega$ and the two components of $\sigma'$). So
locally every metric can be brought to a chosen \emph{fiducial metric}\index{fiducial metric}
$\hat g$; this is the slice. Two caveats are postponed: globally this fails on worldsheets of
higher genus, where a finite number of parameters (moduli) survive, and even locally a
residual symmetry, the conformal group, remains \source{Tong §5.1.1, ll.~6354--6378}.

Define the Faddeev--Popov determinant\index{Faddeev--Popov determinant} $\Delta_{FP}[g]$ by
\begin{equation}\label{eq:ghosts:fpdef}
  \Delta_{FP}[g]^{-1}=\int\mathcal{D}\zeta\;\delta\bigl(g-\hat g^{\zeta}\bigr),
\end{equation}
the exact analogue of $\Delta^{-1}=\int\dd\theta\,\delta(y^\theta)$
\source{Tong §5.1.1, ll.~6379--6392}. The measure $\mathcal D\zeta$ is assumed invariant
under the group, $\mathcal D\zeta=\mathcal D(\zeta'\zeta)$, like $\dd\theta$ above.

\begin{lemma}\label{lem:ghosts:inv}
$\Delta_{FP}$ is gauge invariant: $\Delta_{FP}[g^\zeta]=\Delta_{FP}[g]$.
\end{lemma}
\begin{proof}
$\Delta_{FP}[g^\zeta]^{-1}=\int\mathcal D\zeta'\,\delta(g^\zeta-\hat g^{\zeta'})
=\int\mathcal D\zeta'\,\delta(g-\hat g^{\zeta^{-1}\zeta'})
=\int\mathcal D\zeta''\,\delta(g-\hat g^{\zeta''})$, where the second step acts with
$\zeta^{-1}$ on both arguments of the delta functional and the third uses the invariance of
the measure with $\zeta''=\zeta^{-1}\zeta'$ \source{Tong §5.1.1, ll.~6406--6419}.
\end{proof}

Insert $1=\Delta_{FP}[g]\int\mathcal D\zeta\,\delta(g-\hat g^\zeta)$ into
\cref{eq:ghosts:Z}. The delta functional performs the integral over $g$ and sets
$g=\hat g^\zeta$:
\begin{align}
  Z[\hat g]&=\frac{1}{\mathrm{Vol}}\int\mathcal D\zeta\,\mathcal DX\;
     \Delta_{FP}[\hat g^\zeta]\;\ee^{-S_{\mathrm{Poly}}[X,\hat g^\zeta]}\notag\\
  &=\frac{1}{\mathrm{Vol}}\int\mathcal D\zeta\,\mathcal DX\;
     \Delta_{FP}[\hat g]\;\ee^{-S_{\mathrm{Poly}}[X,\hat g]}
  =\int\mathcal DX\;\Delta_{FP}[\hat g]\;\ee^{-S_{\mathrm{Poly}}[X,\hat g]} .
  \label{eq:ghosts:Zfixed}
\end{align}
The second line uses \cref{lem:ghosts:inv} and the gauge invariance of the action (together
with a relabelling of the integration variable $X$); after that nothing depends on $\zeta$
and $\int\mathcal D\zeta=\mathrm{Vol}$ cancels the prefactor
\source{Tong §5.1.1, ll.~6421--6480}.

\begin{pitfallbox}[Common pitfall: the hidden assumption]
The step ``the action and the measure are gauge invariant'' is exactly where an anomaly would
enter. The combined measure $\mathcal DX\,\Delta_{FP}$ is invariant under diffeomorphisms, but
under a Weyl transformation it changes by the anomaly of \cref{eq:ghosts:weylanomaly}, with $c$
the \emph{total} central charge. The result $Z[\hat g]$ is independent of the arbitrary choice
of $\hat g$, as a gauge-fixed answer must be, only if that total vanishes. The rest of this
chapter computes the total.
\end{pitfallbox}

\section{From a determinant to ghost fields}\label{sec:ghosts:ghosts}

\subsection{Evaluating the determinant}

Near the identity, a gauge transformation is an infinitesimal Weyl rescaling $\omega(\sigma)$
and an infinitesimal diffeomorphism $\delta\sigma^\alpha=v^\alpha(\sigma)$. The metric changes
by
\begin{equation}\label{eq:ghosts:deltag}
  \delta\hat g_{\alpha\beta}=2\omega\,\hat g_{\alpha\beta}+\nabla_\alpha v_\beta+\nabla_\beta v_\alpha ,
\end{equation}
By \cref{lem:ghosts:inv} we need $\Delta_{FP}$ only at $g=\hat g$. There, assuming the slice
meets each orbit once, the delta functional in \cref{eq:ghosts:fpdef} is supported at the
identity transformation, so it suffices to expand $\hat g^\zeta$ to first order, and the
integral over the group becomes an integral over its Lie algebra, $\Delta_{FP}[\hat g]^{-1}=\int\mathcal D\omega\,\mathcal Dv\;
\delta(2\omega\hat g_{\alpha\beta}+\nabla_\alpha v_\beta+\nabla_\beta v_\alpha)$
\source{Tong §5.1.2, ll.~6482--6500}. A delta functional restricts a symmetric tensor field
to vanish, so its Fourier representation needs a symmetric tensor field $\beta^{\alpha\beta}$
as the conjugate variable:
\[
  \Delta_{FP}[\hat g]^{-1}=\int\mathcal D\omega\,\mathcal Dv\,\mathcal D\beta\;
  \exp\Bigl(2\pi\ii\int\dd^2\sigma\sqrt{\hat g}\;\beta^{\alpha\beta}
  \bigl[2\omega\hat g_{\alpha\beta}+\nabla_\alpha v_\beta+\nabla_\beta v_\alpha\bigr]\Bigr).
\]
The field $\omega$ appears without derivatives and linearly. Integrating over it gives
another delta functional, which sets $\beta^{\alpha\beta}\hat g_{\alpha\beta}=0$: from now on
$\beta^{\alpha\beta}$ is symmetric and \emph{traceless}. Using the symmetry of $\beta$ to
combine the two remaining terms,
\begin{equation}\label{eq:ghosts:fpinv}
  \Delta_{FP}[\hat g]^{-1}=\int\mathcal Dv\,\mathcal D\beta\;
  \exp\Bigl(4\pi\ii\int\dd^2\sigma\sqrt{\hat g}\;\beta^{\alpha\beta}\nabla_\alpha v_\beta\Bigr)
\end{equation}
\source{Tong §5.1.2, ll.~6501--6527}. The counting is consistent: in two dimensions a
symmetric traceless tensor has $3-1=2$ components and so does the vector $v_\beta$. The
operator in the exponent maps vectors to symmetric traceless tensors and can be treated as a
square matrix, which is necessary for a determinant to make sense.

\subsection{Commuting versus anticommuting variables}

\Cref{eq:ghosts:fpinv} is a Gaussian integral over commuting variables, which produces an
\emph{inverse} determinant. We want the determinant itself. The tool is integration over
Grassmann (anticommuting) variables\index{Grassmann variables}. For a single pair
$\theta,\bar\theta$ with $\theta^2=\bar\theta^2=0$, $\theta\bar\theta=-\bar\theta\theta$, and the
rules $\int\dd\theta\,1=0$, $\int\dd\theta\,\theta=1$, one has
\[
  \int\dd\bar\theta\,\dd\theta\;\ee^{-\bar\theta a\theta}
  =\int\dd\bar\theta\,\dd\theta\,\bigl(1-a\,\bar\theta\theta\bigr)
  =a\int\dd\bar\theta\,\dd\theta\;\theta\bar\theta = a ,
\]
whereas for a complex commuting variable
$\int\dd\bar z\,\dd z\,\ee^{-\bar z a z}\propto 1/a$. For $n$ pairs and a matrix $M$ the same
computation gives
\begin{equation}\label{eq:ghosts:berezin}
  \int\prod_i\dd\bar\theta_i\,\dd\theta_i\;\ee^{-\bar\theta M\theta}=\det M,
  \qquad
  \int\prod_i\dd\bar z_i\,\dd z_i\;\ee^{-\bar z M z}\propto\frac{1}{\det M},
\end{equation}
because expanding the exponential to order $n$ and keeping the terms that contain every
$\theta_i$ and every $\bar\theta_i$ once reproduces the signed sum over permutations that
defines the determinant (\cref{exr:ghosts:berezin} asks for the $2\times2$ case). This is
standard material, not taken from a source in this book's library.

So to invert \cref{eq:ghosts:fpinv} we replace the commuting fields by anticommuting fields
with the same index structure,
\[
  \beta_{\alpha\beta}\longrightarrow b_{\alpha\beta},\qquad v^\alpha\longrightarrow c^\alpha ,
\]
and obtain, after a rescaling of the fields that changes only the overall normalization of
$Z$ and after rotating back to Euclidean signature,
\begin{equation}\label{eq:ghosts:Sghost}
  \Delta_{FP}[\hat g]=\int\mathcal Db\,\mathcal Dc\;\ee^{-S_{\mathrm{ghost}}},\qquad
  S_{\mathrm{ghost}}=\frac{1}{2\pi}\int\dd^2\sigma\sqrt{g}\;b_{\alpha\beta}\nabla^\alpha c^\beta
\end{equation}
\source{Tong §5.1.3, ll.~6529--6560}. The fields $b_{\alpha\beta}$ (symmetric, traceless)
and $c^\alpha$ are the Faddeev--Popov ghosts\index{ghosts!Faddeev--Popov}\index{bc ghost system}.
The gauge-fixed partition function \cref{eq:ghosts:Zfixed} is
\begin{equation}\label{eq:ghosts:Zfull}
  Z[\hat g]=\int\mathcal DX\,\mathcal Db\,\mathcal Dc\;
  \exp\bigl(-S_{\mathrm{Poly}}[X,\hat g]-S_{\mathrm{ghost}}[b,c,\hat g]\bigr).
\end{equation}

\begin{physicsbox}[Physical picture: what ghosts are for]
Ghosts are anticommuting fields of integer spin. They violate the spin--statistics theorem,
so they cannot describe physical particles, and they are not meant to: they are a local way
of writing a Jacobian. Their effect is to cancel the unphysical, gauge degrees of freedom of
$X^\mu$, leaving the $D-2$ transverse oscillations found in lightcone gauge, but without
breaking manifest Lorentz invariance \source{Tong §5.1.3, ll.~6566--6574}. Tong's
one-loop partition function is computed in lightcone gauge with $24$ oscillators precisely
to avoid carrying the ghosts along \source{Tong §6.4.2, ll.~8519--8527}; in the covariant
computation the ghost oscillators remove the contribution of two of the $26$ directions
(standard; not checked against a source in this book's library). Do not confuse these ghosts with the negative-norm states of covariant quantization
in \cref{ch:quantum}, which are also called ghosts; the Faddeev--Popov ghosts are the cure,
those states are the disease.
\end{physicsbox}

\subsection{Conformal gauge}

Choose the fiducial metric $\hat g_{\alpha\beta}=\ee^{2\omega}\delta_{\alpha\beta}$ and
complex coordinates $z=\sigma^1+\ii\sigma^2$ as in \cref{ch:cft}. Tracelessness,
$\hat g^{\alpha\beta}b_{\alpha\beta}\propto b_{z\bar z}=0$, leaves the components $b_{zz}$ and
$b_{\bar z\bar z}$. With $\sqrt{\hat g}=\ee^{2\omega}$, $\dd^2\sigma=\tfrac12\dd^2z$ and
$\nabla^z=\hat g^{z\bar z}\nabla_{\bar z}=2\ee^{-2\omega}\nabla_{\bar z}$, the factors of
$\ee^{2\omega}$ cancel:
\[
  S_{\mathrm{ghost}}=\frac{1}{2\pi}\int\dd^2z\;
  \bigl(b_{zz}\nabla_{\bar z}c^z+b_{\bar z\bar z}\nabla_z c^{\bar z}\bigr).
\]
The covariant derivative is $\nabla_{\bar z}c^z=\partial_{\bar z}c^z+\Gamma^z_{\bar z\alpha}c^\alpha$,
and $\Gamma^z_{\bar z\alpha}=\tfrac12 g^{z\bar z}(\partial_{\bar z}g_{\alpha\bar z}
+\partial_\alpha g_{\bar z\bar z}-\partial_{\bar z}g_{\bar z\alpha})=0$: the first and third
terms cancel and $g_{\bar z\bar z}=0$ in conformal gauge. Writing
$b=b_{zz}$, $c=c^z$, $\bar b=b_{\bar z\bar z}$, $\bar c=c^{\bar z}$, $\partial=\partial_z$ and
$\bar\partial=\partial_{\bar z}$,
\begin{equation}\label{eq:ghosts:Sbc}
  S_{\mathrm{ghost}}=\frac{1}{2\pi}\int\dd^2z\;\bigl(b\,\bar\partial c+\bar b\,\partial\bar c\bigr)
\end{equation}
\source{Tong §5.1.3--5.2, ll.~6578--6640}. Two observations. The conformal factor $\omega$
has dropped out, so the ghost action is Weyl invariant with $b_{\alpha\beta}$ and $c^\alpha$
inert; this holds for this index placement only, since raising or lowering an index
introduces the metric. And the equations of motion $\bar\partial b=\bar\partial c=0$,
$\partial\bar b=\partial\bar c=0$ say that $b,c$ are holomorphic and $\bar b,\bar c$
antiholomorphic: the ghosts form a free conformal field theory, to which all the tools of
\cref{ch:cft} apply.

\section{The ghost conformal field theory}\label{sec:ghosts:cft}

We treat the holomorphic half; the antiholomorphic half is identical. It costs nothing to be
slightly more general than the string requires, and the generality pays off in
\cref{sec:ghosts:super}. Consider anticommuting fields $b$ and $c$ with action
$\frac{1}{2\pi}\int\dd^2z\,b\,\bar\partial c$, and suppose $b$ has conformal weight $\lambda$ and
$c$ has weight $1-\lambda$, so that $b\,\bar\partial c$ has weight $(1,1)$ and the action is
conformally invariant. The reparametrization ghosts have $\lambda=2$, because $b_{zz}$ has two
lower holomorphic indices and $c^z$ one upper index.

\subsection{Propagator}

As in \cref{ch:cft}, the operator product follows from the statement that the path integral
of a total functional derivative vanishes:
\[
  0=\int\mathcal Db\,\mathcal Dc\;\frac{\delta}{\delta b(\sigma)}
  \Bigl[\ee^{-S_{\mathrm{ghost}}}\,b(\sigma')\Bigr]
  =\int\mathcal Db\,\mathcal Dc\;\ee^{-S_{\mathrm{ghost}}}
  \Bigl[-\frac{1}{2\pi}\bar\partial c(\sigma)\,b(\sigma')+\delta(\sigma-\sigma')\Bigr],
\]
that is, $\bar\partial c(\sigma)\,b(\sigma')=2\pi\,\delta(\sigma-\sigma')$ inside correlation
functions. Since $\bar\partial(1/z)=2\pi\,\delta(z,\bar z)$, this integrates to
\begin{equation}\label{eq:ghosts:bcope}
  c(z)\,b(w)\sim\frac{1}{z-w},\qquad b(z)\,c(w)\sim\frac{1}{z-w},
\end{equation}
where $\sim$ means ``equal up to terms regular as $z\to w$''. The second relation follows
from the first and anticommutativity: $b(z)c(w)=-c(w)b(z)\sim-1/(w-z)=1/(z-w)$. The products
$b(z)b(w)$ and $c(z)c(w)$ are regular \source{Tong §5.2, ll.~6674--6720}.\index{operator product expansion!ghosts}

\subsection{Stress tensor and conformal weights}

The stress tensor is
\begin{equation}\label{eq:ghosts:T}
  T(z)=-\lambda\,{:}b\,\partial c{:}(z)+(1-\lambda)\,{:}\partial b\,c{:}(z) .
\end{equation}
For $\lambda=2$ this is $T=-2\,{:}b\,\partial c{:}-{:}\partial b\,c{:}
=2\,{:}\partial c\,b{:}+{:}c\,\partial b{:}$, which is the result Tong obtains by varying the
covariant action \cref{eq:ghosts:Sghost} with respect to the metric
\source{Tong §5.2, ll.~6650--6672, 6721--6726}. For general $\lambda$ we do not derive
\cref{eq:ghosts:T} from an action; we justify it by checking that it assigns the intended
weights. A primary field $\mathcal O$ of weight $h$ satisfies
$T(z)\mathcal O(w)\sim h\,\mathcal O(w)/(z-w)^2+\partial\mathcal O(w)/(z-w)$
(\cref{ch:cft}).

\emph{The field $b$.} Contract the $c$ in each term of $T(z)$ with $b(w)$, using
$\partial c(z)\,b(w)\sim\partial_z(z-w)^{-1}=-(z-w)^{-2}$:
\[
  T(z)\,b(w)\sim-\lambda\,b(z)\cdot\frac{-1}{(z-w)^2}+(1-\lambda)\,\frac{\partial b(z)}{z-w}
  =\frac{\lambda\,b(w)}{(z-w)^2}+\frac{\bigl[\lambda+(1-\lambda)\bigr]\partial b(w)}{z-w},
\]
after the Taylor expansion $b(z)=b(w)+(z-w)\partial b(w)+\dots$. So $b$ is primary of weight
$\lambda$.

\emph{The field $c$.} Now the $b$ in each term of $T(z)$ is contracted with $c(w)$, and
$c(w)$ must first be moved past one anticommuting field, which costs a sign:
${:}b\,\partial c{:}(z)\,c(w)\to-\partial c(z)/(z-w)$ and
${:}\partial b\,c{:}(z)\,c(w)\to-\partial_z(z-w)^{-1}c(z)=c(z)/(z-w)^2$. Hence
\[
  T(z)\,c(w)\sim\frac{\lambda\,\partial c(z)}{z-w}+\frac{(1-\lambda)\,c(z)}{(z-w)^2}
  =\frac{(1-\lambda)\,c(w)}{(z-w)^2}+\frac{\partial c(w)}{z-w},
\]
so $c$ is primary of weight $1-\lambda$. For $\lambda=2$: $h_b=2$ and $h_c=-1$
\source{Tong §5.2, ll.~6727--6769}. These are the weights dictated by the index structure
$b_{zz}$, $c^z$: a field with $n$ lower $z$ indices has weight $n$, and an upper index counts
as $-1$. A negative weight is a sign that the theory is not unitary, which is what we need to
evade the bound $c>0$.

\subsection{The central charge}

The central charge is read off from the most singular term of
$T(z)T(w)\sim\frac{c/2}{(z-w)^4}+\frac{2T(w)}{(z-w)^2}+\frac{\partial T(w)}{z-w}$
\source{Tong §4.4, ll.~4895--4915}. The $(z-w)^{-4}$ term comes from contracting both fields
of $T(z)$ with both fields of $T(w)$. There are four products of terms, and in each exactly
one pairing is possible ($b$ with $c$). In every case the field being moved passes two
anticommuting fields, so the reordering sign is $+$. With $G=(z-w)^{-1}$ and
\[
  \partial_wG=\frac{1}{(z-w)^2},\qquad\partial_zG=\frac{-1}{(z-w)^2},\qquad
  \partial_z\partial_wG=\frac{-2}{(z-w)^3},
\]
the four double contractions are as follows. In the column ``pairings'' the first factor is
the contraction of the $b$-type field at $z$ with the $c$-type field at $w$, and the second
that of the $c$-type field at $z$ with the $b$-type field at $w$; each is a derivative of
\cref{eq:ghosts:bcope}.
\begin{center}
\renewcommand{\arraystretch}{1.25}
\begin{tabular}{@{}llll@{}}
\toprule
product & pairings & value $\times(z-w)^4$ & at $\lambda=2$\\
\midrule
$\lambda^2\;{:}b\,\partial c{:}(z)\,{:}b\,\partial c{:}(w)$ &
  $\partial_wG\cdot\partial_zG$ & $-\lambda^2$ & $-4$\\
$-\lambda(1-\lambda)\;{:}b\,\partial c{:}(z)\,{:}\partial b\,c{:}(w)$ &
  $G\cdot\partial_z\partial_wG$ & $+2\lambda(1-\lambda)$ & $-4$\\
$-\lambda(1-\lambda)\;{:}\partial b\,c{:}(z)\,{:}b\,\partial c{:}(w)$ &
  $\partial_z\partial_wG\cdot G$ & $+2\lambda(1-\lambda)$ & $-4$\\
$(1-\lambda)^2\;{:}\partial b\,c{:}(z)\,{:}\partial b\,c{:}(w)$ &
  $\partial_zG\cdot\partial_wG$ & $-(1-\lambda)^2$ & $-1$\\
\bottomrule
\end{tabular}
\end{center}
The sum is $c/2=-\lambda^2+4\lambda(1-\lambda)-(1-\lambda)^2=-6\lambda^2+6\lambda-1$, that is
\begin{equation}\label{eq:ghosts:cbc}
  \boxed{\,c_{bc}(\lambda)=-2\,(6\lambda^2-6\lambda+1)=1-3\,(2\lambda-1)^2\,}
\end{equation}
(the two forms agree: $1-3(4\lambda^2-4\lambda+1)=-12\lambda^2+12\lambda-2$; the table and
this identity were also confirmed by a symbolic check in sympy). For the reparametrization
ghosts the last column adds up to $-13$, so\index{central charge!of the bc ghosts}
\begin{equation}\label{eq:ghosts:minus26}
  T(z)\,T(w)\sim\frac{-13}{(z-w)^4}+\frac{2T(w)}{(z-w)^2}+\frac{\partial T(w)}{z-w},
  \qquad c_{bc}(2)=-26,
\end{equation}
in agreement with \source{Tong §5.2, ll.~6770--6839}, who carries out the $\lambda=2$ case
including the less singular terms. The general-$\lambda$ formula is standard; here it is
derived on the page rather than quoted.

\begin{example}[The formula at the weights that matter]\label{ex:ghosts:table}
\Cref{eq:ghosts:cbc} is a downward parabola in $\lambda$, symmetric under
$\lambda\leftrightarrow1-\lambda$ (exchange of the roles of $b$ and $c$), with maximum $c=1$
at $\lambda=\tfrac12$ (\cref{fig:ghosts:parabola}).
\begin{center}
\begin{tabular}{@{}ccccl@{}}
\toprule
$\lambda$ & $1-\lambda$ & $(2\lambda-1)^2$ & $c_{bc}(\lambda)$ & system\\
\midrule
$1/2$ & $1/2$ & $0$ & $+1$ & a complex free fermion $=$ two real fermions\\
$1$ & $0$ & $1$ & $-2$ & \\
$3/2$ & $-1/2$ & $4$ & $-11$ & (the commuting version has $+11$, see \cref{sec:ghosts:super})\\
$2$ & $-1$ & $9$ & $-26$ & reparametrization ghosts\\
\bottomrule
\end{tabular}
\end{center}
The first row is a useful by-product. At $\lambda=\tfrac12$ the fields $b$ and $c$ both have
weight $\tfrac12$ and can be combined into two real (Majorana) fermions
$\psi_{1,2}=(b\pm c)/\sqrt2$ up to a factor of $\ii$. So one real free fermion has
\[
  c_\psi=\tfrac12 ,
\]
the value used for the superstring below.\index{central charge!of a free fermion}
\end{example}

\begin{figure}[ht]
\centering
\begin{tikzpicture}[xscale=2.6,yscale=0.14]
  \draw[-{Stealth}] (-1.25,0) -- (2.4,0) node[right] {$\lambda$};
  \draw[-{Stealth}] (0,-29) -- (0,5) node[above] {$c_{bc}(\lambda)$};
  \draw[thick,domain=-1.02:2.02,samples=60] plot (\x,{1-3*(2*\x-1)*(2*\x-1)});
  \foreach \x/\y/\lab/\pos in {0.5/1/{$(\tfrac12,1)$}/above right, 1.5/-11/{$(\tfrac32,-11)$}/right,
     2/-26/{$(2,-26)$}/right, 1/-2/{$(1,-2)$}/right}
     {\fill (\x,\y) circle[x radius=0.025, y radius=0.46]; \node[\pos,font=\small] at (\x,\y) {\lab};}
  \fill (-1,-26) circle[x radius=0.025, y radius=0.46];
  \node[left,font=\small] at (-1,-26) {$(-1,-26)$};
  \draw[dashed] (-1,-26) -- (2,-26);
  \foreach \y in {-26,-11} {\draw (-0.03,\y) -- (0.03,\y);}
  \node[left,font=\small] at (-0.03,-11) {$-11$};
\end{tikzpicture}
\caption{The central charge $c_{bc}(\lambda)=1-3(2\lambda-1)^2$ of an anticommuting
first-order system whose fields have weights $\lambda$ and $1-\lambda$. The points
$\lambda=2$ and $\lambda=-1$ describe the same system with the names of $b$ and $c$
exchanged. The commuting system of the same weights has the opposite central charge.}
\label{fig:ghosts:parabola}
\end{figure}

\section{The critical central charge}\label{sec:ghosts:critical}

Collect the pieces. After gauge fixing, the worldsheet theory in \cref{eq:ghosts:Zfull} consists
of a ``matter'' conformal field theory (so far, $D$ free scalars) and the ghosts. Central
charges of decoupled theories add, because their stress tensors add and have no singular
product with each other. The Weyl anomaly of the full theory, \cref{eq:ghosts:weylanomaly}, is
proportional to
\begin{equation}\label{eq:ghosts:ctot}
  c_{\mathrm{tot}}=c_{\mathrm{matter}}+c_{bc}(2)=c_{\mathrm{matter}}-26 ,
\end{equation}
and the gauge symmetry we used to fix the metric survives quantization only if
$c_{\mathrm{tot}}=0$:

\begin{proposition}[Critical central charge \TL]\label{prop:ghosts:26}
The gauge-fixed bosonic string is Weyl invariant if and only if the matter conformal field
theory has $c_{\mathrm{matter}}=\tilde c_{\mathrm{matter}}=26$. If the matter theory consists
of $D$ free scalars, each with $c=1$, this requires $D=26$
\source{Tong §5.3, ll.~6840--6851}.\index{critical dimension}\index{critical central charge}
\end{proposition}

This agrees with the lightcone result of \cref{ch:quantum}, but the logic is different and
more informative. There, $D=26$ was needed to keep a global symmetry (Lorentz invariance)
that the gauge choice had hidden. Here, Lorentz invariance is manifest throughout and
$c=26$ is needed to keep the gauge symmetry itself.

\begin{physicsbox}[Physical picture: a budget, not a dimension]
Nothing in the argument required the matter theory to be free scalars. \emph{Any} conformal
field theory with $c=\tilde c=26$ gives a consistent gauge-fixed string, and each describes a
different background in which the string propagates. ``Critical dimension'' is thus a
misnomer for ``critical central charge'': only for special conformal field theories can $c$
be read as a number of spacetime dimensions. To describe strings in four-dimensional
Minkowski space one takes four free scalars and an ``internal'' conformal field theory with
$c=26-4=22$, which may or may not have a geometric interpretation
\source{Tong §5.3, ll.~6852--6867}. When this book compactifies on a circle
(\cref{ch:circle}), a torus (\cref{ch:narain}) or a K3 surface (\cref{ch:stringsk3}), it is
choosing the internal theory. T-duality will turn out to be the statement that two
different-looking geometries give the \emph{same} internal conformal field theory.
\end{physicsbox}

\begin{remark}[Time]\label{rem:ghosts:time}
One direction must be timelike, i.e.\ have a wrong-sign kinetic term. The safe construction
keeps $X^0$ as a free scalar and lets the remaining $c=25$ be an arbitrary conformal field
theory, which excludes backgrounds that evolve in time. The worldsheet description of
time-dependent backgrounds has unresolved technical obstacles, and string cosmology is
usually studied through the low-energy effective action of \cref{ch:backgrounds} instead
\source{Tong §5.3, ll.~6868--6877}. This caveat applies to the cosmological discussion of
\cref{ch:cosmology}.
\end{remark}

\begin{remark}[Non-critical strings]\label{rem:ghosts:liouville}
If $c_{\mathrm{matter}}=D\neq26$ one may give up Weyl invariance from the start. The
conformal factor $\omega$ of the metric then does not decouple: integrating the anomaly
equation with $c=D-26$ shows that $\omega$ acquires a kinetic term proportional to
$(26-D)$ and, if a worldsheet cosmological constant $\mu$ is present, an exponential
potential $\mu\,\ee^{2\omega}$. It becomes a new dynamical field, the Liouville
field\index{Liouville field} \source{Tong §5.3.2, ll.~6892--6987}. We shall not use
non-critical strings; the remark shows that ``$D\neq26$'' does not simply fail but turns into
a different and harder theory.
\end{remark}

\subsection{Physical states without BRST}

With the ghosts present the space of states is larger, not smaller, and one needs a criterion
for physical states. The systematic answer is BRST quantization\index{BRST quantization},
which neither Tong nor this book develops \source{Tong §5.4, ll.~6988--7001}. The following
shortcut is enough for our purposes. By the state--operator map of \cref{ch:cft}, a state
corresponds to a local operator $\mathcal O(z,\bar z)$ of the matter theory. An operator at a
point is not diffeomorphism invariant; its integral over the worldsheet is:
\begin{equation}\label{eq:ghosts:vertex}
  V\sim\int\dd^2z\;\mathcal O(z,\bar z).
\end{equation}
Under a rescaling the measure $\dd^2z$ has weight $(-1,-1)$, so Weyl invariance of $V$
requires $\mathcal O$ to be a primary operator of weight $(h,\tilde h)=(1,1)$. Such $V$ are
the \emph{vertex operators}\index{vertex operator} \source{Tong §5.4, ll.~7020--7059}.

\begin{example}[Masses from weights]\label{ex:ghosts:masses}
The operator ${:}\ee^{\ii p\cdot X}{:}$ has weight $h=\tilde h=\ap p^2/4$ (\cref{ch:cft}).
Setting $h=1$ gives $p^2=4/\ap$, i.e.\ $M^2=-p^2=-4/\ap$: the tachyon. The operator
${:}\ee^{\ii p\cdot X}\partial X^\mu\bar\partial X^\nu{:}\,\zeta_{\mu\nu}$ has
$h=\tilde h=1+\ap p^2/4$, so $h=1$ gives $M^2=0$; demanding that it be primary (no
$(z-w)^{-3}$ term in its product with $T$) imposes the transversality conditions
$p^\mu\zeta_{\mu\nu}=p^\nu\zeta_{\mu\nu}=0$ \source{Tong §5.4.1, ll.~7060--7116}. Both
masses agree with the book's mass formula $M^2=\frac{2}{\ap}(N+\tilde N-2)$ at
$N=\tilde N=0$ and $N=\tilde N=1$. The intercept ``$-2$'' of that formula, which in
\cref{ch:quantum} came from a regularized sum over zero-point energies, comes here from the
requirement of weight $(1,1)$.
\end{example}

Ghosts can also be ignored in tree-level scattering: on the sphere they decouple from the
matter fields, whereas on surfaces of higher genus the moduli of the metric couple the two,
in analogy with the field-theory statement that ghosts matter only in loops
\source{Tong §6.1.1, ll.~7444--7458}.

\section{The superstring in brief}\label{sec:ghosts:super}

This book works mostly with the bosonic fields $G$, $B$ and the dilaton, which are common to
all string theories, so a full treatment of the superstring is not needed. But the K3 part of
the book (\cref{ch:stringsk3,ch:k3t2,ch:ellgenus}) takes place in superstring theory, and
the student needs to know what the words mean. Tong's notes, our main source, deliberately
give ``nothing more than a list of facts'' about the superstring
\source{Tong §2.5, ll.~2958--2961}; we do the same, and mark what goes beyond that source.
Everything in this section is \TL.

\subsection{Worldsheet fermions and \texorpdfstring{$D=10$}{D = 10}}

In the Neveu--Schwarz--Ramond formulation the worldsheet theory contains, besides the scalars
$X^\mu$, anticommuting fields $\psi^\mu$ of weight $\tfrac12$ (and/or their antiholomorphic
partners $\tilde\psi^\mu$), related to $X^\mu$ by a worldsheet supersymmetry\index{superstring}
\index{worldsheet supersymmetry}. This supersymmetry is a further \emph{gauge} symmetry on the
worldsheet, so gauge fixing it produces further ghosts, the $\beta\gamma$
system\index{beta-gamma ghost system@$\beta\gamma$ ghost system}, with central charge $+11$.
The matter theory must therefore have
\[
  c_{\mathrm{matter}}=26-11=15 .
\]
Supersymmetry pairs each boson ($c=1$) with a fermion ($c=\tfrac12$; compare
\cref{ex:ghosts:table}), so $D$ flat dimensions contribute $\tfrac32D$, and
\begin{equation}\label{eq:ghosts:D10}
  \tfrac32\,D=15\quad\Longrightarrow\quad D=10
\end{equation}
\source{Tong §5.3.1, ll.~6878--6891}.\index{critical dimension!of the superstring}

Where does $+11$ come from? Tong states the number without derivation. It follows from
\cref{sec:ghosts:cft} with two changes. (i) The gauge parameter of a local supersymmetry is a
spinor, a ``square root'' of the vector $c^z$ of weight $-1$; its ghost $\gamma$ has weight
$-\tfrac12$, and the conjugate field $\beta$ has weight $\lambda=\tfrac32$. (ii) Ghosts have
the opposite statistics to the gauge parameter, so $\beta,\gamma$ \emph{commute}. For
commuting fields with the same action $\frac{1}{2\pi}\int\dd^2z\,\beta\bar\partial\gamma$ the
derivation of \cref{eq:ghosts:bcope} still gives $\gamma(z)\beta(w)\sim1/(z-w)$, but now
$\beta(z)\gamma(w)=+\gamma(w)\beta(z)\sim1/(w-z)=-1/(z-w)$. In the computation of the weights
this sign replaces the reordering sign of the anticommuting case, so the weights are still
$(\lambda,1-\lambda)$. In each of the four double contractions of the table in
\cref{sec:ghosts:cft} exactly one of the two pairings is of the type $\beta(z)\gamma(w)$, so
every entry changes sign:
\begin{equation}\label{eq:ghosts:cbg}
  c_{\beta\gamma}(\lambda)=-c_{bc}(\lambda)=3(2\lambda-1)^2-1,\qquad
  c_{\beta\gamma}\bigl(\tfrac32\bigr)=3\cdot4-1=11 .
\end{equation}
The identification of the weights in (i) is standard but is not checked against a source in
this book's library; the sign argument in (ii) is derived here.

\begin{example}[The budget for $\KT$]\label{ex:ghosts:budget}
Take the superstring on $\R^{1,3}\times T^2\times K3$ (\cref{ch:k3t2}). The ghosts contribute
$-26+11=-15$. The matter side is
\begin{center}
\begin{tabular}{@{}lccc@{}}
\toprule
factor & real dimensions & $c$ per dimension & $c$\\
\midrule
$\R^{1,3}$ & $4$ & $1+\tfrac12$ & $6$\\
$T^2$ & $2$ & $1+\tfrac12$ & $3$\\
K3 & $4$ & $1+\tfrac12$ & $6$\\
\midrule
total & $10$ & & $15$\\
\bottomrule
\end{tabular}
\end{center}
and $15-15=0$. The entry for K3 deserves a comment. A K3 surface is curved, and its sigma
model is an interacting theory, not four free bosons and four free fermions. The value $c=6$
is nevertheless correct, because the anomaly budget leaves exactly $15-9=6$ for the internal
theory once $\R^{1,3}\times T^2$ is accounted for, and because a Ricci-flat target preserves
conformal invariance to the order discussed in \cref{ch:backgrounds}. The free-field count
$4\cdot1+4\cdot\tfrac12$ is a mnemonic (it is exact for the orbifold $T^4/\Z_2$ of
\cref{ch:kummer}). The same number $c=6$ labels the $N=4$ superconformal algebra whose
characters organize the elliptic genus in \cref{ch:ellgenus}.
\end{example}

\subsection{The five theories}

The bosonic string is unique, but there are discrete choices in adding fermions to the
worldsheet. The main one is whether fermions are added to both the left- and right-moving
sectors or to one sector only \source{Tong §2.5, ll.~2969--2985}. In all cases $D=10$,
there is no tachyon, and the massless bosonic fields include the metric $G_{\mu\nu}$, the
two-form $B_{\mu\nu}$ and the dilaton $\Phi$ \source{Tong §2.5, ll.~2958--2968}, which is why
the results of \cref{ch:backgrounds,ch:tduality,ch:buscher} carry over.
\index{type IIA string}\index{type IIB string}\index{heterotic string}\index{type I string}

\begin{center}
\renewcommand{\arraystretch}{1.2}
\begin{tabular}{@{}p{26mm}p{30mm}p{18mm}p{60mm}@{}}
\toprule
theory & worldsheet fermions & super\-charges & additional massless bosons\\
\midrule
type IIA & both sectors & $32$ & Ramond--Ramond one-form $C_1$ and three-form $C_3$\\
type IIB & both sectors & $32$ & Ramond--Ramond $C_0$, $C_2$, $C_4$; the field strength of
  $C_4$ is self-dual, $F_5=\star F_5$\\
heterotic $SO(32)$ & one sector & $16$ & $SO(32)$ gauge field\\
heterotic $E_8\times E_8$ & one sector & $16$ & $E_8\times E_8$ gauge field\\
type I & (open and closed strings) & & mentioned by Tong only in passing\\
\bottomrule
\end{tabular}
\end{center}
\source{Tong §2.5, ll.~2981--3008}. Three comments connect this table to later chapters.

\emph{Heterotic strings and lattices.} In a heterotic string one sector is that of the
superstring, with matter central charge $15$, and the other is that of the bosonic string,
with matter central charge $26$. Ten dimensions use $10$ of the $26$, and the remaining
$26-10=16$ units are carried by an internal theory that produces the gauge group; both
$SO(32)$ and $E_8\times E_8$ have rank $16$. The consistent choices are governed by even
unimodular lattices of rank $16$, the subject of \cref{ch:lattices}, and on a torus by the
Narain lattices of \cref{ch:narain}. (This paragraph is standard; Tong gives only the list of
gauge groups, and the counting is not checked against a source in this book's library.)

\emph{T-duality.} The superstring theories are not individually invariant under
$R\to\ap/R$; they are exchanged in pairs. Type IIA on a circle of radius $R$ is equivalent to
type IIB on a circle of radius $\ap/R$, and the two heterotic strings are likewise exchanged
\source{Tong §8.3.3, ll.~11828--11835}. This is consistent with the D-brane content: type IIA
has stable D$p$-branes with $p$ even, type IIB with $p$ odd, and T-duality along a circle
changes $p$ by one (\cref{ch:dbranes}).

\emph{K3.} Aspinwall's lectures, the main source for \cref{ch:stringsk3}, state their goal as
understanding ``the type II string and the heterotic string compactified on a K3 surface and
what such models are dual to'', and remark that the heterotic string appears to be
intrinsically tied to the geometry of K3 \source{Asp §1, ll.~112--190}. The student who knows
the table above knows enough to follow that discussion.

\begin{pitfallbox}[Common pitfall: two meanings of ``supersymmetry'']
The supersymmetry used in \cref{eq:ghosts:D10} lives on the worldsheet and is a gauge
redundancy. The ``$32$ supercharges'' of the table refer to spacetime supersymmetry of the
ten-dimensional theory, which is a physical symmetry. The second does not follow from the
first by inspection; it requires a projection on the spectrum that this book does not
describe (standard; not covered by the sources in this book's library).
\end{pitfallbox}

\begin{historybox}
The analysis of the string path integral described in this chapter is due to Polyakov,
``Quantum geometry of bosonic strings'', Phys.\ Lett.\ B \textbf{103}, 207 (1981), a paper of
four pages, followed by ``Quantum geometry of fermionic strings'', Phys.\ Lett.\ B
\textbf{103}, 211 (1981), for the superstring. Polyakov's main interest was the non-critical
string of \cref{rem:ghosts:liouville}, not $D=26$
\source{Tong §5.1, ll.~6312--6315; §5.3.2, ll.~6892--6896}.
\end{historybox}

\section{What the machine has checked}\label{sec:ghosts:lean}

Two files of the library \lean{StringTheoryFormalization} concern this chapter. Both model
central charges as rational numbers and prove arithmetic facts about them. In the language of
this book they are \emph{arithmetic shadows}: they check the last line of the argument, not
the argument. No path integral, no Grassmann field, no operator product and no Virasoro
algebra is defined in either file.

\begin{leanbox}[In Lean: \texttt{UseCases/CriticalDimension.lean}]
The two central-charge formulas \cref{eq:ghosts:cbc,eq:ghosts:cbg} are \emph{definitions}:
\begin{leancode}
def fermionicGhostCharge (lam : ℚ) : ℚ := 1 - 3 * (2 * lam - 1) ^ 2
def bosonicGhostCharge (lam : ℚ) : ℚ := 3 * (2 * lam - 1) ^ 2 - 1
\end{leancode}
The theorems proved about them, in the namespace
\lean{StringTheory.UseCases.CriticalDimension}, are:
\begin{leancode}
theorem bosonic_eq_neg_fermionic (lam : ℚ) :
    bosonicGhostCharge lam = - fermionicGhostCharge lam := by ...

theorem bc_ghost_charge_eq : fermionicGhostCharge 2 = -26 := by ...

theorem betagamma_ghost_charge_eq : bosonicGhostCharge (3 / 2) = 11 := by ...

theorem bosonic_string_critical_dimension :
    (26 : ℚ) * 1 + fermionicGhostCharge 2 = 0 := by ...

theorem superstring_critical_dimension :
    (10 : ℚ) * 1 + 10 * (1 / 2) + fermionicGhostCharge 2
      + bosonicGhostCharge (3 / 2) = 0 := by ...

theorem bosonic_dimension_unique (D : ℚ)
    (h : D * 1 + fermionicGhostCharge 2 = 0) :
    D = 26 := by ...
\end{leancode}
\textbf{Plain reading.} The polynomial $3(2\lambda-1)^2-1$ is the negative of
$1-3(2\lambda-1)^2$ for every rational $\lambda$; the first polynomial equals $11$ at
$\lambda=\tfrac32$ and the second equals $-26$ at $\lambda=2$; $26-26=0$;
$10+5-26+11=0$; and the only rational $D$ with $D-26=0$ is $26$. The numerals $1$ and
$\tfrac12$ multiplying $26$ and $10$ stand for the central charges of a free boson and a free
fermion, but in the statements they are just numerals.

\textbf{Proof idea.} Unfold the definitions; \lean{ring} proves the polynomial identity,
\lean{norm\_num} evaluates at a fixed rational number, and \lean{linarith} solves the linear
equation in \lean{bosonic\_dimension\_unique}.

\textbf{What is not proved.} That a $bc$ system of weight $\lambda$ has central charge
\lean{fermionicGhostCharge}~$\lambda$ (our derivation of \cref{eq:ghosts:cbc} is on paper
only); that the Faddeev--Popov ghosts have $\lambda=2$; that the Weyl anomaly is proportional
to the total central charge; or that a consistent conformal field theory with $c=26$ exists.
There is also no uniqueness statement for the superstring in the file
(\cref{exr:ghosts:leanunique}).
\end{leanbox}

\begin{leanbox}[In Lean: \texttt{Frontier/CentralCharge.lean}]
In the namespace \lean{StringTheory.Frontier}:
\begin{leancode}
def centralChargeBoson : ℚ := 1
def centralChargeFermion : ℚ := 1/2
def centralChargeK3 : ℚ :=
  4 * centralChargeBoson + 4 * centralChargeFermion

theorem central_charge_k3_eq_six : centralChargeK3 = 6 := by ...
\end{leancode}
\textbf{Plain reading.} $4\cdot1+4\cdot\tfrac12=6$ in $\Q$. This is the K3 row of
\cref{ex:ghosts:budget}, with the free-field mnemonic taken as the definition.

The same file contains three theorems whose names promise more than their statements
deliver, as the file's own documentation says explicitly:
\begin{leancode}
theorem virasoro_ope_central_term (z w : ℂ) (h : z ≠ w) :
    (centralChargeK3 / 2 : ℚ) = 3 := by ...

theorem virasoro_algebra_commutator (c : ℚ) (m n : ℤ) (L : ℤ → ℤ) :
    c / 12 * (m^3 - m) = c * m * (m^2 - 1) / 12 := by ...

theorem central_charge_from_worldsheet_action :
    ∃ (c : ℚ), c = centralChargeK3 ∧ c = 6 := by ...
\end{leancode}
The first does not use $z$, $w$ or $h$ and states that $6/2=3$. The second does not use $n$
or $L$ and is the identity $\frac{c}{12}(m^3-m)=\frac{c\,m(m^2-1)}{12}$, valid for every $c$;
it says nothing about commutators $[L_m,L_n]$. The third is witnessed by $c=6$ and involves
no action. They are recorded here so that a reader who meets these names in the catalogue
(\cref{ch:atlas}) is not misled. One sentence of that file's documentation should also be read
with care: it describes the bookkeeping as ``$26=20+6$'', which mixes the bosonic budget $26$
with the superstring value $c=6$. The consistent superstring bookkeeping is $15=9+6$, as in
\cref{ex:ghosts:budget}. This affects a comment only, not a theorem.
\end{leanbox}

All declarations above depend only on the axioms \lean{propext}, \lean{Classical.choice} and
\lean{Quot.sound}.

\begin{center}
\renewcommand{\arraystretch}{1.2}
\begin{tabular}{@{}p{70mm}c>{\raggedright\arraybackslash\footnotesize}p{62mm}@{}}
\toprule
statement & tier & evidence\\
\midrule
$1-3(2\lambda-1)^2$ equals $-26$ at $\lambda=2$; its negative equals $11$ at $\tfrac32$ & \TA &
  \lean{bc\_ghost\_charge\_eq}, \lean{betagamma\_ghost\_charge\_eq}\\
$26-26=0$; $10+5-26+11=0$; $D-26=0\Rightarrow D=26$ & \TA &
  \lean{bosonic\_string\_critical\_dimension}, \lean{superstring\_critical\_dimension},
  \lean{bosonic\_dimension\_unique}\\
$4\cdot1+4\cdot\frac12=6$ & \TA & \lean{central\_charge\_k3\_eq\_six}\\
$c_{bc}(\lambda)=1-3(2\lambda-1)^2$ for a $bc$ system of weight $\lambda$ & \TL &
  derived in \cref{sec:ghosts:cft}; sympy check; $\lambda=2$ in Tong §5.2\\
Gauge fixing yields $bc$ ghosts with $\lambda=2$; Weyl invariance iff $c_{\mathrm{matter}}=26$ &
  \TL & Tong §5.1--5.3\\
$\beta\gamma$ ghosts have $c=11$; superstring has $D=10$ & \TL & Tong §5.3.1\\
Identification of the rational numbers in the Lean files with central charges of worldsheet
theories & --- & a modelling statement, never Tier A\\
\bottomrule
\end{tabular}
\end{center}

\begin{remark}[What a faithful formalization would need]\label{rem:ghosts:faithful}
To make \cref{prop:ghosts:26} a theorem one would need, at the least: a definition of a
(chiral) conformal field theory or vertex operator algebra with a stress tensor and its
central charge; the construction of the $bc$ system as such an object for each $\lambda$; a
proof that its central charge is \cref{eq:ghosts:cbc}, which is a finite computation with
formal power series once the definitions exist; and additivity of $c$ under tensor products.
None of these is available in Mathlib at the time of writing, to our knowledge. The
path-integral steps of \cref{sec:ghosts:fp,sec:ghosts:ghosts} are formal manipulations even
on paper and are further still from formalization. \Cref{ch:open} returns to this.
\end{remark}

\begin{summarybox}
\begin{itemize}[leftmargin=*,itemsep=2pt]
\item Diffeomorphisms and Weyl transformations are gauge redundancies of the Polyakov path
  integral. Dividing by their volume is done by the Faddeev--Popov method: choose a fiducial
  metric $\hat g$ and pay the Jacobian $\Delta_{FP}[\hat g]$.
\item $\Delta_{FP}$ is the determinant of the operator taking a vector $v$ to the traceless
  part of $\nabla_\alpha v_\beta+\nabla_\beta v_\alpha$. Written as a Grassmann integral it
  becomes the ghost action $\frac{1}{2\pi}\int\dd^2z\,(b\bar\partial c+\bar b\partial\bar c)$.
\item The ghosts form a conformal field theory with $h_b=2$, $h_c=-1$ and $c=-26$. For general
  weight, $c_{bc}(\lambda)=1-3(2\lambda-1)^2$; the commuting system has the opposite sign.
\item The Weyl anomaly is proportional to the total central charge. Consistency requires
  $c_{\mathrm{matter}}=26$ for the bosonic string and $15$ for the superstring; with free
  fields, $D=26$ and $D=10$. These are budgets for a conformal field theory, not only counts
  of dimensions; a K3 surface uses $c=6$.
\item Physical states correspond to matter primaries of weight $(1,1)$.
\item There are five superstring theories in ten dimensions. Type IIA and IIB are exchanged
  by T-duality, as are the two heterotic strings.
\item Lean checks the arithmetic of the final step (values of two polynomials, three linear
  identities, one uniqueness statement). Everything connecting that arithmetic to the
  worldsheet is Tier L.
\end{itemize}
\end{summarybox}

\section*{Exercises}
\addcontentsline{toc}{section}{Exercises}

\begin{exercise}[$\star$]\label{exr:ghosts:plane3}
Repeat \cref{ex:ghosts:plane} in $\R^3$, declaring the direction of $\vec x$ to be pure
gauge, so that the orbits are spheres and $\mathrm{Vol}=4\pi$. Compute
$Z=\frac{1}{4\pi}\int\dd^3x\,\ee^{-|\vec x|^2}$ directly and with the slice ``positive
$z$-axis'', and identify the Faddeev--Popov Jacobian. (Answer: $\Delta=r^2$,
$Z=\sqrt\pi/4$.)
\end{exercise}

\begin{exercise}[$\star$]\label{exr:ghosts:berezin}
Verify \cref{eq:ghosts:berezin} for $n=2$: expand
$\exp(-\sum_{ij}\bar\theta_iM_{ij}\theta_j)$ to second order and show that the coefficient of
$\theta_1\bar\theta_1\theta_2\bar\theta_2$ is $M_{11}M_{22}-M_{12}M_{21}$.
\end{exercise}

\begin{exercise}[$\star$]\label{exr:ghosts:count}
In $n$ worldsheet dimensions, count the components of a symmetric traceless tensor and of a
vector. Show that they agree only for $n=2$. What does this say about the attempt to gauge
away the whole metric of a membrane worldvolume ($n=3$) with diffeomorphisms and Weyl
transformations?
\end{exercise}

\begin{exercise}[$\star\star$]\label{exr:ghosts:lesssing}
For $\lambda=2$, compute the $(z-w)^{-2}$ term of $T(z)T(w)$ from the single contractions and
show that it equals $2T(w)$. You will need to Taylor expand fields at $z$ around $w$ and to
keep track of one reordering sign per term.
\end{exercise}

\begin{exercise}[$\star\star$]\label{exr:ghosts:symmetry}
Show directly from \cref{eq:ghosts:T} that exchanging the names $b\leftrightarrow c$ together
with $\lambda\to1-\lambda$ maps $T$ to itself (use ${:}c\,\partial b{:}=-{:}\partial b\,c{:}$),
and conclude that $c_{bc}(\lambda)=c_{bc}(1-\lambda)$. Then state and prove in Lean:
\begin{leancode}
theorem fermionicGhostCharge_symm (lam : ℚ) :
    fermionicGhostCharge lam = fermionicGhostCharge (1 - lam) := by
  sorry
\end{leancode}
(Hint: \lean{unfold fermionicGhostCharge; ring}.)
\end{exercise}

\begin{exercise}[$\star\star$]\label{exr:ghosts:leanunique}
The Lean file proves uniqueness of $D=26$ but not of $D=10$. State and prove the missing
theorem: for $D:\Q$, if
\lean{D * 1 + D * (1 / 2) + fermionicGhostCharge 2 + bosonicGhostCharge (3 / 2) = 0} then
$D=10$. Which two existing theorems should you rewrite with before calling \lean{linarith}?
\end{exercise}

\begin{exercise}[$\star\star$]\label{exr:ghosts:fermion}
Prove in Lean that \lean{fermionicGhostCharge (1 / 2) = 1}. Explain in two sentences why
this supports the value $c=\tfrac12$ for one real fermion, which the second file records as
the constant \lean{centralChargeFermion}. Is your explanation part of what Lean checked?
\end{exercise}

\begin{exercise}[$\star\star$]\label{exr:ghosts:internal}
(a) For the bosonic string on $\R^{1,d-1}$ times an internal theory, give the internal central
charge as a function of $d$. (b) Do the same for the superstring, and find the internal central
charge for $d=4$ and $d=6$. (c) A Calabi--Yau threefold has real dimension $6$. Check that its
free-field count of $c$ matches (b) for $d=4$, and that K3 matches $d=6$.
\end{exercise}

\begin{exercise}[$\star\star\star$]\label{exr:ghosts:liouville}
Starting from the anomaly equation
$\frac{1}{Z}\frac{\partial Z}{\partial\omega}=\frac{c}{24\pi}\sqrt g\,(R-2\nabla^2\omega)$
with $\mu=0$ and $c=D-26$ (\cref{rem:ghosts:liouville}), integrate in $\omega$ to show that
\[
  Z[\ee^{2\omega}g]=Z[g]\,\exp\Bigl(\frac{c}{24\pi}\int\dd^2\sigma\sqrt g\;
  \bigl(g^{\alpha\beta}\partial_\alpha\omega\,\partial_\beta\omega+R\,\omega\bigr)\Bigr).
\]
For which sign of
$D-26$ does the Liouville field have a kinetic term of the conventional sign? Compare with
\source{Tong §5.3.2, ll.~6940--6985}.
\end{exercise}

\section*{Further reading}
\addcontentsline{toc}{section}{Further reading}

\begin{itemize}[leftmargin=*,itemsep=2pt]
\item \source{Tong §5.1, ll.~6300--6613}: the Faddeev--Popov procedure for the string, in the
  form followed here. \source{Tong §5.2, ll.~6614--6839}: the ghost conformal field theory and
  the full $TT$ product for $\lambda=2$.
\item \source{Tong §5.3, ll.~6840--6987}: critical central charge, the superstring count and
  non-critical strings. \source{Tong §5.4, ll.~6988--7116}: vertex operators.
\item \source{Tong §2.5, ll.~2948--3008} and \source{Tong §8.3.3, ll.~11828--11835}: the list
  of superstring theories and their behaviour under T-duality.
\item \source{Tong §6.1.1, ll.~7329--7412}: the sum over worldsheet topologies and the string
  coupling, which complete the definition of the path integral begun here.
\item \source{Asp §1, ll.~112--190}: why type II and heterotic strings on K3 are the natural
  setting for the second half of this book.
\item BRST quantization and the general first-order system are treated in Polchinski's
  textbook, to which Tong refers \source{Tong §5.4, ll.~6997--7001}; that book is not in
  this book's library, and the Lean file cites its equations (2.7.20) and (10.4.10) for the
  two ghost-charge formulas without our having checked them against it.
\end{itemize}
